\section{Introduction}
\label{sec:intro}
The increasing frequency and severity of extreme weather events, driven by anthropogenic climate change, pose unprecedented challenges to societal infrastructure and human safety \cite{IPCC2021}. Effective early-warning mechanisms for hazards such as floods, severe heatwaves, and their compounding effects are critical for proactive disaster management. While modern machine learning and deep learning architectures have demonstrated substantial theoretical capacity for environmental prediction, their transition from offline scientific studies to realtime operational systems remains fraught with methodological bottlenecks \cite{Zscheischler2020}. 

Two fundamental challenges constrain the deployment of contemporary predictive systems. First, evaluating these models on highly autocorrelated environmental time-series data using randomized cross-validation often introduces look-ahead bias, producing artificially inflated metrics that collapse during live deployment. Second, advanced temporal sequence models, such as Long Short-Term Memory (LSTM) networks or Gated Recurrent Units (GRU), achieve high predictive accuracy by maintaining deep historical states \cite{Kratzert2018}. However, orchestrating the continuous retrieval, buffering, and processing of these historical sequences introduces severe latency and complexity constraints in live, stateless operational environments.

To address these limitations, this study proposes ClimateGuardian AI, a hybrid predictive framework integrating traditional machine learning and deep learning methodologies. The system is designed to provide robust, multi-hazard risk assessment utilizing a 45-feature hydro-climatic dataset while strictly enforcing a leakage-aware chronological evaluation protocol. By deliberately distinguishing between theoretical offline sequence benchmarks and deployment-ready stateless algorithms, the framework successfully bridges the gap between historical predictive evaluation and live Open-Meteo operational inference. 
